\section{Regularización y generalización}

\subsection{El problema del sobreajuste}

Al entrenar redes neuronales, un desafío central es el \textbf{sobreajuste} (Overfitting): el modelo tiene un buen desempeño en los datos de entrenamiento, pero un desempeño deficiente en datos de prueba que nunca ha visto. Esto significa que el modelo "memorizó" el ruido y los patrones particulares de los datos de entrenamiento, sin aprender reglas verdaderamente generalizables.

\begin{importantbox}{Subajuste, ajuste adecuado y sobreajuste}
\begin{itemize}
\item \textbf{Subajuste} (Underfitting): el modelo es demasiado simple para capturar los patrones fundamentales de los datos y tiene un desempeño deficiente tanto en el conjunto de entrenamiento como en el de prueba
\item \textbf{Ajuste adecuado} (Good Fit): el modelo aprende de manera apropiada los patrones centrales de los datos y tiene buen desempeño tanto en el conjunto de entrenamiento como en el de prueba
\item \textbf{Sobreajuste} (Overfitting): el modelo es demasiado complejo y se ajusta al ruido de los datos de entrenamiento; tiene buen desempeño en el conjunto de entrenamiento pero deficiente en el de prueba
\end{itemize}
\end{importantbox}

\subsection{Técnicas de regularización}

Para mitigar el sobreajuste, es necesario usar técnicas de \textbf{regularización} (Regularization). Amini presentó dos métodos principales:

\textbf{1. Dropout}

La idea central de Dropout es sumamente sencilla: durante el entrenamiento, con probabilidad $p$ se "apaga" (se pone en cero) de manera aleatoria una parte de las neuronas.

\begin{itemize}
\item En la propagación hacia adelante de cada mini-batch, se elige de manera aleatoria alrededor de $p\%$ de las neuronas para que su salida sea cero
\item Esto obliga a la red a no depender de ninguna neurona individual, sino a aprender representaciones de características más distribuidas y robustas
\item Durante la prueba no se usa Dropout, pero es necesario escalar la salida de manera proporcional
\end{itemize}

\begin{lstlisting}[caption={Implementación en código de Dropout}]
# TensorFlow
tf.keras.layers.Dropout(rate=0.5)

# PyTorch
nn.Dropout(p=0.5)
\end{lstlisting}

\textbf{2. Early Stopping}

Early Stopping monitorea el desempeño del modelo en el conjunto de validación: cuando la pérdida de validación deja de disminuir (o incluso empieza a aumentar), se detiene el entrenamiento. Esto aprovecha una observación clave: el sobreajuste suele ocurrir en las etapas finales del entrenamiento, momento en el que el modelo empieza a "memorizar" los datos de entrenamiento en lugar de aprender reglas generalizables.

\begin{warningbox}{La regularización no es opcional}
En la práctica, una red neuronal profunda sin regularización casi con certeza se sobreajustará, en especial cuando el número de parámetros del modelo es mucho mayor que el número de ejemplos de entrenamiento (algo muy común en el aprendizaje profundo actual). Dropout y Early Stopping deben ser prácticas estándar, no elementos adicionales opcionales.
\end{warningbox}

\subsection{Resumen de la sección}

El sobreajuste es un desafío central en el entrenamiento del aprendizaje profundo. Dropout obliga a la red a aprender características robustas apagando neuronas de manera aleatoria, y Early Stopping detiene el entrenamiento en el momento adecuado monitoreando la pérdida en el conjunto de validación. La combinación de ambos puede mejorar de manera significativa la capacidad de generalización del modelo.
